\input format.tex

%\usepackage{arydshln}

\begin{document}

%%%
\begin{spacing}{1.25}

\setlength{\arrayrulewidth}{0.5pt}
\color{gray2}
\vspace*{8mm}
\fontsize{9.3pt}{9.8pt}\selectfont
\begin{tctabularx}{tabularx={p{15.4cm}<{\centering}p{0pt}}}
\\[-6.769pt]
\cellcolor{topcolor} \raisebox{1pt}{\color{gray0}\xiaosihao\lantxh\bf 益生菌调理方案} \\
\end{tctabularx}

{
\lantxh
\vspace*{-4.25mm}
\noindent
\begin{longtable}{|p{2.25cm}|p{12.732cm}|}
\hline
\hspace*{5mm}\parbox[c]{\hsize}{\vskip10pt\raisebox{-1.5ex}{\includegraphics[scale=1.25]{prebiotics_fangan.pdf}}\\{\xiaosihao\bf 益生菌}\vskip10pt} &
\parbox[c]{\hsize}{\vskip10pt\wuhao {建议每天服用适量益生菌OTC产品（如双歧杆菌活菌制剂、乳杆菌活菌制剂等），可以调节肠道PH值，增加肠道内益生菌的数量，帮助抑制肠道有害菌的生长，有利于肠道菌群平衡，提高肠道菌群多样性。饭后1小时食用效果更佳。}\vskip10pt}\\
\hline
\end{longtable}
}

\vspace*{8mm}
\fontsize{9.3pt}{9.8pt}\selectfont
\begin{tctabularx}{tabularx={p{15.4cm}<{\centering}p{0pt}}}
\\[-6.769pt]
\cellcolor{topcolor} \raisebox{1pt}{\color{gray0}\xiaosihao\lantxh\bf 益生元调理方案} \\
\end{tctabularx}

{
\lantxh
\vspace*{-4.25mm}
\noindent
\begin{longtable}{|p{2.25cm}|p{12.732cm}|}
\hline
\hspace*{5mm}\parbox[c]{\hsize}{\vskip10pt\raisebox{-1.5ex}{\includegraphics[scale=1.25]{probiotics_fangan.pdf}}\\{\xiaosihao\bf 益生元}\vskip10pt} &
\parbox[c]{\hsize}{\vskip10pt\wuhao {建议每天服用适量低聚糖类产品，如低聚果糖、低聚半乳糖、低聚异麦芽糖、低聚甘露糖，促进体内双歧杆菌、乳杆菌等益生菌的生长，提高菌群多样性。早晚空腹服用效果更佳。如患炎症性肠病、慢性腹泻，不宜服用菊粉，建议选择低聚半乳糖、低聚果糖。}\vskip10pt}\\
\hline
\end{longtable}
}

\end{spacing}
\end{document}
